% !TEX root = ../../../main.tex

\section{非功能性测试}

非功能性测试主要用于验证 HearBridge 系统在性能、稳定性、安全性、易用性和可维护性等方面的表现。

与黑盒功能测试不同，
非功能性测试不只关注某个功能是否能够完成，
还关注系统在连续访问、异常输入、权限控制和后续维护等场景下是否具备基本可用性。

本节结合系统实际运行情况，
分别对性能、稳定性、安全性、易用性和可维护性进行测试与分析。

\subsection{性能测试}

性能测试主要用于验证 HearBridge 后端核心接口在一定并发访问条件下的响应效率和稳定性。

本次性能测试使用 Apache JMeter 5.6.3 完成。

测试对象选择 Spring Boot 后端中较常用的管理端与资源查询接口，
包括管理员登录、当前管理员信息查询、手势分类分页查询、手势分类列表查询、手势资源分页查询和手势资源列表查询。

上述接口覆盖了后台管理端登录鉴权和移动端资源浏览中的高频访问场景，
能够反映系统后端基础业务接口的响应能力。

表~\ref{tab:chapter07-jmeter-test-config} 给出了 JMeter 性能测试配置。

\begin{longtable}{L{0.24\textwidth} L{0.28\textwidth} L{0.38\textwidth}}
  \caption{JMeter 性能测试配置表}
  \label{tab:chapter07-jmeter-test-config}\\
  \toprule
  配置项 & 配置内容 & 说明 \\
  \midrule
  测试工具 & Apache JMeter 5.6.3 & 用于构造并发请求并生成测试报告 \\
  测试目标 & Spring Boot 后端核心接口 & 包括登录、鉴权信息查询、分类查询和资源查询等接口 \\
  测试地址 & \texttt{http://127.0.0.1:8080} & 本地后端服务地址 \\
  20 并发场景 & 20 个线程，Ramp-Up 20 秒，循环 10 次 & 用于验证基础并发访问能力 \\
  50 并发场景 & 50 个线程，Ramp-Up 50 秒，循环 10 次 & 用于验证中等并发访问能力 \\
  结果统计方式 & JMeter HTML Dashboard & 统计错误率、平均响应时间、百分位响应时间和吞吐量 \\
  \bottomrule
\end{longtable}

性能测试过程中，
JMeter 首先调用管理员登录接口获取 token，
随后在同一测试流程中请求当前管理员信息、手势分类和手势资源相关接口。

这种测试方式能够同时覆盖公开查询接口和需要鉴权的管理端接口，
比单纯测试公开接口更接近后台管理端的实际访问场景。

图~\ref{fig:chapter07-jmeter-dashboard-50u} 展示了 50 并发场景下的 JMeter 总体统计结果。

\WhuAutoFigure{0.88\textwidth}{0.48\textheight}{figures/chapter07/jmeter-dashboard-50u.png}{JMeter 50 并发总体统计截图}{fig:chapter07-jmeter-dashboard-50u}

表~\ref{tab:chapter07-performance-result} 给出了 20 并发和 50 并发两组测试结果。

\begin{longtable}{L{0.14\textwidth} R{0.10\textwidth} R{0.10\textwidth} R{0.13\textwidth} R{0.12\textwidth} R{0.12\textwidth} R{0.12\textwidth} R{0.12\textwidth}}
  \caption{后端核心接口性能测试结果表}
  \label{tab:chapter07-performance-result}\\
  \toprule
  测试场景 & 并发数 & 请求数 & 平均响应/ms & 中位数/ms & 90\%响应/ms & 吞吐量/s & 错误率 \\
  \midrule
  基础并发 & 20 & 1020 & 16.11 & 5.00 & 18.00 & 53.11 & 0.00\% \\
  中等并发 & 50 & 2550 & 10.50 & 4.00 & 8.00 & 51.78 & 0.00\% \\
  \bottomrule
\end{longtable}

由表~\ref{tab:chapter07-performance-result} 可知，
在 20 并发和 50 并发测试条件下，
系统核心接口错误率均为 0.00\%。

在 20 并发场景下，
总请求数为 1020，
平均响应时间为 16.11 ms，
90\% 响应时间为 18.00 ms，
吞吐量为 53.11 transactions/s。

在 50 并发场景下，
总请求数为 2550，
平均响应时间为 10.50 ms，
90\% 响应时间为 8.00 ms，
吞吐量为 51.78 transactions/s。

两组测试中，
系统均未出现请求失败，
说明后端核心查询接口在本地测试环境下具有较好的响应性能。

图~\ref{fig:chapter07-jmeter-response-time} 展示了性能测试过程中的响应时间变化情况。

\WhuAutoFigure{0.88\textwidth}{0.48\textheight}{figures/chapter07/jmeter-response-times-over-time-50u.png}{JMeter 响应时间趋势图}{fig:chapter07-jmeter-response-time}

从图~\ref{fig:chapter07-jmeter-response-time} 可以看出，
大部分请求响应时间保持在较低水平。

测试过程中虽然存在少量响应时间波动，
但整体没有出现持续升高或大量请求阻塞的情况。

图~\ref{fig:chapter07-jmeter-percentiles} 展示了响应时间百分位变化情况。

\WhuAutoFigure{0.88\textwidth}{0.48\textheight}{figures/chapter07/jmeter-percentiles-over-time-50u.png}{JMeter 响应时间百分位变化图}{fig:chapter07-jmeter-percentiles}

从图~\ref{fig:chapter07-jmeter-percentiles} 可以看出，
多数请求的响应时间分布较集中，
高百分位响应时间没有出现持续异常升高。

结合总体统计结果可知，
HearBridge 后端核心接口能够满足毕业设计演示环境下的基础并发访问需求。

需要说明的是，
本次性能测试主要针对常用后端业务接口展开。

模型训练、句子视频识别和实时图像识别等计算密集型任务，
其耗时与视频长度、模型规模、关键点检测速度和硬件性能密切相关，
不适合与普通业务查询接口混合进行同一组 JMeter 压力测试。

因此，
本节性能测试结论主要用于说明后端核心业务接口的并发访问能力，
不代表模型推理类任务的极限性能。

\subsection{稳定性与可靠性测试}

稳定性与可靠性测试主要用于验证系统在连续操作、重复请求和异常场景下是否能够保持正常运行。

测试过程中，
分别对移动端、管理端、后端接口和 Python 识别服务进行了多轮操作。

移动端测试包括重复进入训练页面、反复打开和关闭相机 overlay、重复发起实时识别、重复选择句子视频并执行识别等操作。

管理端测试包括多次刷新分类列表、资源列表、样本列表、模型版本列表，
以及重复执行查询、打开弹窗、取消操作和返回页面等操作。

后端接口测试包括多次请求分类分页、资源分页、登录接口和模型版本接口，
观察接口是否存在异常中断、响应超时或服务崩溃情况。

Python 识别服务测试包括服务启动、模型加载、实时识别连接、样本采集、模型重载和句子视频识别等流程。

表~\ref{tab:chapter07-stability-test} 给出了稳定性与可靠性测试结果。

\begin{longtable}{L{0.22\textwidth} L{0.34\textwidth} L{0.34\textwidth}}
  \caption{稳定性与可靠性测试结果表}
  \label{tab:chapter07-stability-test}\\
  \toprule
  测试对象 & 测试内容 & 测试结果 \\
  \midrule
  移动端训练页面 & 多次进入训练页、播放数字人动作、打开和关闭相机 overlay & 页面能够正常响应，未出现严重卡死或崩溃 \\
  移动端实时识别 & 重复开启和停止实时识别流程 & WebSocket 连接和识别状态能够正常切换 \\
  移动端句子识别 & 多次选择视频、快速识别和执行 AI 忠实自然化 & 页面能够展示结果或错误提示，异常情况下未崩溃 \\
  后台管理端 & 多次刷新列表、打开弹窗、提交表单和取消操作 & 页面状态正常，表格和弹窗能够正确刷新 \\
  后端接口 & 连续请求核心查询接口和认证接口 & 接口能够持续返回结果，JMeter 测试错误率为 0.00\% \\
  Python 识别服务 & 启动服务、加载模型、执行识别和重载模型 & 服务能够按预期响应请求，异常输入下能够返回错误信息 \\
  \bottomrule
\end{longtable}

测试结果表明，
HearBridge 系统在常规使用场景下能够保持稳定运行。

在异常输入或外部服务异常情况下，
系统能够通过错误提示、异常返回或状态重置等方式进行处理，
未出现影响整体运行的严重问题。

\subsection{安全性测试}

安全性测试主要验证系统对未授权访问、错误登录、无效 token 和异常请求的处理能力。

HearBridge 后端通过统一认证拦截器区分公开接口、移动端用户接口和管理端接口。

公开查询接口允许直接访问；
移动端接口需要 App 用户 token；
管理端接口需要管理员 token。

当请求未携带 token 或 token 无效时，
系统应拒绝访问并返回未授权结果。

表~\ref{tab:chapter07-security-test} 给出了安全性测试结果。

\begin{longtable}{L{0.22\textwidth} L{0.34\textwidth} L{0.34\textwidth}}
  \caption{安全性测试结果表}
  \label{tab:chapter07-security-test}\\
  \toprule
  测试内容 & 操作方式 & 预期与实际结果 \\
  \midrule
  管理员错误登录 & 使用错误用户名或错误密码访问管理端登录接口 & 系统拒绝登录，并返回用户名或密码错误提示 \\
  管理端未登录访问 & 未携带管理员 token 访问管理端受保护接口 & 系统返回未授权结果，前端跳转登录页 \\
  无效 token 访问 & 使用空 token、过期 token 或错误 token 请求受保护接口 & 系统拒绝访问，并返回登录失效或未授权提示 \\
  退出登录后访问 & 管理员退出登录后继续访问受保护接口 & 系统不再允许访问，需要重新登录 \\
  表单参数校验 & 新增分类或资源时提交空编码、空名称或非法编码 & 页面或后端返回校验提示，不保存非法数据 \\
  文件上传校验 & 上传不符合格式要求的 SiGML 文件 & 页面拒绝上传，并提示文件格式或大小不符合要求 \\
  \bottomrule
\end{longtable}

由表~\ref{tab:chapter07-security-test} 可知，
系统能够对常见非法访问和异常输入进行拦截。

管理端路由守卫能够阻止未登录用户直接访问业务页面；
后端认证拦截器能够对受保护接口进行 token 校验；
前端表单校验和后端异常处理能够减少非法数据进入系统。

因此，
系统在毕业设计阶段具备基本的访问控制和输入校验能力。

\subsection{易用性测试}

易用性测试主要从普通用户和管理员两个角度观察系统操作流程是否清晰、页面反馈是否明确。

移动端普通用户侧主要关注首页入口、训练流程、数字人展示、实时识别和句子视频识别等功能是否易于理解。

后台管理端管理员侧主要关注分类、资源、样本、训练和模型版本等模块的操作入口是否清晰，
表格、弹窗和提示信息是否便于使用。

表~\ref{tab:chapter07-usability-test} 给出了易用性测试结果。

\begin{longtable}{L{0.22\textwidth} L{0.34\textwidth} L{0.34\textwidth}}
  \caption{易用性测试结果表}
  \label{tab:chapter07-usability-test}\\
  \toprule
  测试对象 & 测试内容 & 测试结果 \\
  \midrule
  移动端底部导航 & 首页、训练、我的三个 Tab 切换 & 导航结构清晰，用户能够快速进入主要功能 \\
  手势训练页面 & 数字人动作播放、相机打开、实时识别和采集模式 & 功能入口集中，识别状态和模型版本信息展示明确 \\
  句子视频识别页面 & 选择视频、快速识别、AI 忠实自然化和词段详情展示 & 操作流程较清晰，结果展示包含 Gloss、中文、英文和置信说明 \\
  管理端侧边栏 & 分类、资源、样本、训练、模型版本菜单 & 功能划分明确，符合管理员资源维护流程 \\
  管理端表格 & 分页、筛选、刷新和操作按钮 & 表格内容较完整，操作反馈较明确 \\
  管理端弹窗 & 新增、编辑、删除、质量标记和发布确认 & 操作前有确认提示，提交后有成功或失败反馈 \\
  \bottomrule
\end{longtable}

测试结果表明，
HearBridge 系统主要页面入口清晰，
移动端和管理端均能围绕各自使用场景完成操作。

移动端更侧重训练与识别体验，
管理端更侧重资源维护与模型管理流程。

系统在关键操作处提供了按钮状态、加载状态、错误提示和确认弹窗，
能够降低误操作带来的影响。

\subsection{可维护性与扩展性测试}

可维护性与扩展性测试主要结合系统工程结构进行分析。

HearBridge 系统采用多端分离结构，
HarmonyOS 移动端、Vue 管理端、Spring Boot 后端和 Python 识别服务分别承担不同职责。

这种结构使页面展示、业务管理、数据存储和模型识别能够分开维护，
有利于后续功能扩展和问题定位。

在后端部分，
系统按照 Controller、Service、Mapper、Entity、DTO 等结构组织代码，
并针对 Controller、Service 和 Interceptor 编写了自动化测试。

在管理端部分，
系统按照路由页面、业务 API、类型定义和通用组件组织代码，
HearBridge 相关功能集中在分类管理、资源管理、样本管理、训练管理和模型版本管理模块中。

在 Python 识别服务部分，
系统将实时识别、样本采集、模型训练、模型运行时管理和句子视频识别拆分为不同模块，
便于后续替换模型或扩展识别能力。

表~\ref{tab:chapter07-maintainability-test} 给出了可维护性与扩展性分析结果。

\begin{longtable}{L{0.22\textwidth} L{0.34\textwidth} L{0.34\textwidth}}
  \caption{可维护性与扩展性分析表}
  \label{tab:chapter07-maintainability-test}\\
  \toprule
  分析对象 & 当前设计 & 维护与扩展效果 \\
  \midrule
  多端架构 & 移动端、管理端、后端和 Python 服务分离 & 各端职责较清晰，便于独立开发和调试 \\
  后端分层 & Controller、Service、Mapper、DTO 和实体类分层 & 业务逻辑和接口入口分离，便于测试和维护 \\
  管理端模块 & 分类、资源、样本、训练和模型版本模块化组织 & 资源维护链路清晰，后续可继续扩展管理功能 \\
  Python 服务模块 & 实时识别、数据集、训练、模型运行时和句子视频识别分模块组织 & 有利于后续替换模型、扩充词表或增加识别接口 \\
  模型版本机制 & 训练产物登记、版本详情查看和发布状态管理 & 支持后续多个模型版本切换与回滚 \\
  自动化测试 & 后端 108 个测试用例全部通过，并统计覆盖率 & 能够辅助发现核心业务逻辑问题，提高维护可靠性 \\
  \bottomrule
\end{longtable}

总体来看，
HearBridge 系统在当前阶段具备较好的模块划分和基础测试支撑。

后续如果继续扩展系统，
可以在现有结构上增加更多手语资源、更多模型版本、更多训练数据和更完整的连续识别能力。

同时，
还可以继续补充工具类、配置类、异常分支和外部服务异常场景的自动化测试，
进一步提升系统可维护性。

\subsection{非功能性测试结果分析}

综合性能测试、稳定性测试、安全性测试、易用性测试和可维护性分析结果可知，
HearBridge 系统在非功能性方面能够满足毕业设计阶段的使用要求。

在性能方面，
系统后端核心接口在 20 并发和 50 并发条件下错误率均为 0.00\%，
平均响应时间和高百分位响应时间均保持在较低水平，
说明后端基础业务接口具有较好的响应能力。

在稳定性方面，
移动端训练页面、句子视频识别页面、后台管理端页面、后端接口和 Python 识别服务在多轮操作中均能够保持正常响应，
异常情况下能够通过错误提示或状态恢复机制避免页面崩溃。

在安全性方面，
系统能够对错误登录、未授权访问、无效 token、退出登录后的访问和非法表单输入进行基本拦截，
具备基础访问控制能力。

在易用性方面，
移动端围绕训练与识别流程组织页面，
后台管理端围绕资源维护、样本管理、训练任务和模型版本发布组织功能，
整体操作路径较清晰。

在可维护性方面，
系统采用移动端、管理端、后端和 Python 识别服务分离的结构，
并通过后端分层、管理端模块化和模型版本管理机制提高系统后续扩展能力。

因此，
非功能性测试结果表明，
HearBridge 系统不仅能够完成主要业务功能，
也能够在性能、稳定性、安全性和维护性方面支撑当前毕业设计系统的运行与展示。